\documentclass[11pt,a4paper,fontset=none]{ctexart}

\usepackage[margin=1in]{geometry}
\usepackage{hyperref}
\usepackage{xcolor}
\usepackage{enumitem}
\usepackage{longtable}
\usepackage{booktabs}
\usepackage{array}
\usepackage{fancyvrb}
\usepackage{fontspec}
\usepackage{xeCJK}

\hypersetup{
  colorlinks=true,
  linkcolor=blue!60!black,
  urlcolor=blue!60!black
}

\setlist[itemize]{leftmargin=2em}
\setlist[enumerate]{leftmargin=2.2em}
\setlength{\parskip}{0.45em}
\setlength{\parindent}{0pt}

\setmainfont{Times New Roman}
\setsansfont{Arial}
\setmonofont{Menlo}
\setCJKmainfont{Songti SC}
\setCJKsansfont{Hiragino Sans GB}
\setCJKmonofont{Hiragino Sans GB}

\title{Project 3 仓库说明（中文版）}
\author{AIAA 3201 Project 3: Video Object Removal \& Inpainting}
\date{\today}

\begin{document}
\maketitle

\section{文档目的}

这份文档是当前仓库的中文版使用说明，主要回答三类问题：

\begin{itemize}
  \item 仓库里每个文件和目录是做什么用的；
  \item 目前哪些部分已经实现，哪些部分只是为后续实验预留接口；
  \item 因为 GPU 资源在服务器上，接下来你们应该按什么顺序把这套代码迁移到服务器并真正跑起来。
\end{itemize}

如果你们后面忘了某个文件的作用，优先先看这份文档，而不是直接猜。

\section{项目当前总体结构}

当前仓库围绕下面这条主线组织：

\begin{center}
\texttt{动态目标掩码生成 -> 掩码优化 -> 视频补全 -> 定量评测 -> 可视化 -> 报告写作}
\end{center}

整体上已经具备：

\begin{itemize}
  \item Part 1 baseline 的代码骨架与命令入口；
  \item Part 2 的外部模型接入接口；
  \item Part 3 的 failure case 工作区准备工具；
  \item 指标计算、可视化、报告和展示材料的模板。
\end{itemize}

\section{顶层文件说明}

\subsection{\texttt{README.md}}

这是仓库的主入口文档。

主要用途：

\begin{itemize}
  \item 解释整个项目的目标；
  \item 展示目录结构；
  \item 给出环境安装方式；
  \item 给出 Part 1 运行方式；
  \item 给出评测命令；
  \item 给出 Part 2 / Part 3 的推荐工作流。
\end{itemize}

任何第一次打开仓库的人，应该先读它。

\subsection{\texttt{requirements.txt}}

这是 Python 依赖列表。

里面当前列出的核心包包括：

\begin{itemize}
  \item \texttt{opencv-python}：读写视频、光流、掩码处理、图像修复；
  \item \texttt{ultralytics}：YOLOv8-seg；
  \item \texttt{PyYAML}：读取配置文件；
  \item \texttt{scikit-image}：SSIM；
  \item \texttt{numpy / scipy / matplotlib}：基础数值和可视化支持。
\end{itemize}

\subsection{\texttt{.gitignore}}

这个文件是为了避免把大文件误提交到仓库里。

它会忽略：

\begin{itemize}
  \item 原始数据视频；
  \item 抽帧结果；
  \item 实验输出；
  \item 指标结果；
  \item 虚拟环境和缓存文件。
\end{itemize}

因为视频项目的中间产物非常大，这个文件很重要。

\section{配置文件说明}

\subsection{\texttt{configs/project3.yaml}}

这是默认实验配置文件，也是后续最常修改的文件之一。

它控制：

\begin{itemize}
  \item YOLO 模型名称；
  \item 要保留的目标类别；
  \item 光流动态判断阈值；
  \item 掩码膨胀参数；
  \item 时间窗口大小；
  \item Part 2 中 \texttt{SAM 2} 和 \texttt{ProPainter} 的外部路径；
  \item Part 3 的说明性参数。
\end{itemize}

可以把它理解成实验控制面板。

\section{代码入口说明}

\subsection{\texttt{scripts/project3.py}}

这是你们之后最常直接运行的入口脚本。

它封装了命令行接口，目前支持：

\begin{itemize}
  \item \texttt{part1}：运行 Part 1 baseline；
  \item \texttt{eval-mask}：计算 mask 指标；
  \item \texttt{eval-video}：计算恢复结果指标；
  \item \texttt{figures}：生成对比图；
  \item \texttt{part2-prepare}：建立 Part 2 实验工作区；
  \item \texttt{part3-prepare}：建立 Part 3 failure case 工作区。
\end{itemize}

这个脚本存在的意义是让你们不需要手动处理 \texttt{PYTHONPATH}。

\section{\texttt{src/project3/} 代码逐文件说明}

\subsection{\texttt{config.py}}

负责读取配置文件。

除了标准 YAML 读取外，它还额外提供了一个轻量 fallback 解析器。这样做的目的是：即使本地环境暂时没有安装 \texttt{PyYAML}，仓库默认配置也仍然能被读取。

\subsection{\texttt{io\_utils.py}}

这是一些小型文件工具函数。

它负责：

\begin{itemize}
  \item 自动创建目录；
  \item 列出图像文件；
  \item 写 JSON 摘要文件。
\end{itemize}

它本身不做视觉算法，只是供其他模块复用。

\subsection{\texttt{video\_utils.py}}

这是视频层面的底层工具函数。

它负责：

\begin{itemize}
  \item 读取视频并拆成帧；
  \item 把帧写回成视频；
  \item 把一组帧保存成图片序列。
\end{itemize}

所有 Part 1 的读写流程都依赖这个模块。

\subsection{\texttt{part1.py}}

这是当前最核心的实现文件，也是现在真正有算法内容的主体。

它实现了 Part 1 baseline 的完整流程：

\begin{enumerate}
  \item 读取输入视频；
  \item 对每帧运行 YOLOv8-seg；
  \item 保留指定类别的候选目标；
  \item 在候选目标内部用 Lucas-Kanade sparse optical flow 判断该目标是否真的是动态目标；
  \item 用膨胀和连通域清理优化掩码；
  \item 先用邻近帧的时间中值去借背景；
  \item 若仍有空洞，再使用 \texttt{cv2.inpaint} 作为补充；
  \item 保存输入帧、mask、恢复帧、恢复视频和摘要 JSON。
\end{enumerate}

简单说，\texttt{part1.py} 就是你们当前第一版“能跑通、能交差、还能做对比”的方法。

\subsection{\texttt{eval\_metrics.py}}

这个模块负责评测。

它支持两类指标：

\begin{itemize}
  \item mask 指标：\texttt{JM}、\texttt{JR}、precision；
  \item 恢复质量指标：\texttt{PSNR}、\texttt{SSIM}。
\end{itemize}

这个文件主要在有 GT 数据时使用，比如 DAVIS 或你们自己整理出的带真值子集。

\subsection{\texttt{visualization.py}}

这个模块负责把实验结果整理成报告和 PPT 好用的对比图。

当前生成的图会展示：

\begin{itemize}
  \item original frame；
  \item predicted mask；
  \item restored frame。
\end{itemize}

它的意义不是“让代码更好看”，而是让你们写报告时直接有图可用。

\subsection{\texttt{part2.py}}

这个文件不是在本仓库里重写 \texttt{SAM 2} 或 \texttt{ProPainter}，而是做一个“适配层”。

它负责：

\begin{itemize}
  \item 创建 Part 2 的标准输出目录；
  \item 写入 metadata；
  \item 生成推荐的调用命令模板。
\end{itemize}

这样做的目的，是让你们把强模型放在外部仓库里维护，同时保证当前主仓库的组织是统一的。

\subsection{\texttt{part3.py}}

这个模块是给 Part 3 的 failure case 实验准备工作区的。

它负责为某个失败案例创建：

\begin{itemize}
  \item 关键帧目录；
  \item mask 目录；
  \item 编辑后关键帧目录；
  \item 比较结果目录；
  \item 记录这个 failure case 目标的说明 JSON。
\end{itemize}

它的作用是帮助你们把 Part 3 控制在“小范围、有证据、便于展示”的状态，而不是一开始就把 diffusion 扩展做成一个大工程。

\subsection{\texttt{cli.py}}

这是命令行接口的总调度模块。

所有在 \texttt{scripts/project3.py} 里调用的命令，真正都会转到这里来执行。

\subsection{\texttt{\_\_init\_\_.py}}

这是标准的 Python 包标记文件，让 \texttt{project3} 目录可以被 Python 识别成模块。

\section{文档目录说明}

\subsection{\texttt{docs/submission\_checklist.md}}

最终提交前的操作清单。

主要用来确认：

\begin{itemize}
  \item GitHub 仓库是否整理好；
  \item 报告是否完整；
  \item arXiv 是否上传；
  \item \texttt{videos.zip} 是否包含 mandatory 数据集结果；
  \item 图表、链接、引用是否齐全。
\end{itemize}

\subsection{\texttt{docs/presentation\_outline.md}}

这是 8 分钟 progress presentation 的讲稿结构建议。

它的目标不是帮你们写逐字稿，而是提醒你们：在 Week 11/12 的汇报里，最少要覆盖哪些内容。

\subsection{\texttt{docs/pipeline\_notes.md}}

这是一个更偏执行层面的流程说明。

相比 README，它更强调：

\begin{itemize}
  \item Part 1、Part 2、Part 3 应该如何衔接；
  \item 什么时候启动 Part 2；
  \item 什么时候才适合开始 Part 3。
\end{itemize}

\subsection{\texttt{docs/repo\_guide.md}}

英文版仓库说明。

\subsection{\texttt{docs/repo\_guide\_zh.tex}}

也就是你现在在看的这份中文版 LaTeX 文档。

\section{报告相关文件说明}

\subsection{\texttt{report/outline.md}}

这是一个自然语言版的论文结构草稿。

适合你们在真正写论文前，先决定每节写什么。

\subsection{\texttt{report/main.tex}}

这是论文主文档的 LaTeX 骨架。

当前它只是一个轻量起步模板，不是正式的 CVPR 官方模板。

真正交作业时，你们应该把内容迁移到课程要求的 CVPR 模板中。

\subsection{\texttt{report/references.bib}}

这是参考文献占位文件。

后续必须把 project PDF 里列出的 references 全部补进去，并确保正文里都被引用到。

\section{数据和输出目录说明}

\subsection{\texttt{data/raw/}}

放原始视频和 GT 数据。

建议使用：

\begin{itemize}
  \item \texttt{data/raw/wild/}
  \item \texttt{data/raw/bmx/}
  \item \texttt{data/raw/tennis/}
  \item \texttt{data/raw/davis/}
\end{itemize}

\subsection{\texttt{data/frames/}}

如果你们想额外做手动抽帧或中间处理，可以把结果放在这里。

\subsection{\texttt{outputs/part1/}}

存 Part 1 baseline 的输出。

每个实验通常应包含：

\begin{itemize}
  \item input frames；
  \item masks；
  \item restored frames；
  \item restored video；
  \item summary JSON；
  \item comparison figure。
\end{itemize}

\subsection{\texttt{outputs/part2/}}

存 Part 2 的输出，比如：

\begin{itemize}
  \item refined masks；
  \item ProPainter 输入；
  \item ProPainter 输出；
  \item metadata 和推荐命令。
\end{itemize}

\subsection{\texttt{outputs/part3/}}

存 failure case 和关键帧编辑实验的结果。

\subsection{\texttt{metrics/}}

存导出的指标 JSON。

\section{现在已经完成了什么}

目前已经完成的部分：

\begin{itemize}
  \item 仓库结构；
  \item baseline Part 1 代码；
  \item 评测代码；
  \item 可视化代码；
  \item Part 2 / Part 3 工作区准备工具；
  \item 报告和展示用文档骨架。
\end{itemize}

目前还没有在本地真正跑起来的部分：

\begin{itemize}
  \item YOLOv8-seg 实际推理，因为本地尚未安装相关依赖；
  \item SAM 2 / ProPainter，因为它们要在服务器环境和 GPU 上运行；
  \item 完整实验，因为数据还没有放进仓库目录。
\end{itemize}

\section{为什么现在不要继续在本地硬跑}

因为你们的 GPU 资源在服务器上，所以本地当前阶段更适合做：

\begin{itemize}
  \item 仓库整理；
  \item 文档准备；
  \item 代码结构搭建；
  \item 路线设计；
  \item 报告模板和参考文献准备。
\end{itemize}

真正的模型推理、长视频处理、SOTA 模型接入，都应该放到服务器上做。

\section{接下来在服务器上的推荐操作顺序}

\subsection*{Step 1：同步当前仓库到服务器}

把这个仓库 clone 到服务器工作目录，或者直接上传整个项目文件夹。

\subsection*{Step 2：创建环境并安装依赖}

推荐命令：

\begin{Verbatim}[fontsize=\small]
python3 -m venv .venv
source .venv/bin/activate
pip install -r requirements.txt
\end{Verbatim}

如果服务器已有 CUDA 版本的 PyTorch，请优先保证 \texttt{ultralytics}、SAM 2 和 ProPainter 与那套环境兼容。

\subsection*{Step 3：放置 mandatory 数据集}

把视频放到：

\begin{Verbatim}[fontsize=\small]
data/raw/wild/
data/raw/bmx/
data/raw/tennis/
\end{Verbatim}

如果你们想拿更强的量化结果，也把 DAVIS 放好。

\subsection*{Step 4：先跑 Part 1 baseline}

推荐先跑 wild video，再跑 bmx 和 tennis。

示例命令：

\begin{Verbatim}[fontsize=\small]
python3 scripts/project3.py part1 \
  --config configs/project3.yaml \
  --video data/raw/wild/corridor.mp4 \
  --experiment wild_corridor
\end{Verbatim}

先做 Part 1 的原因：

\begin{itemize}
  \item 它最快形成第一个完整闭环；
  \item 它能尽快产出 progress presentation 可用结果；
  \item 它给 Part 2 提供清晰的对比基线。
\end{itemize}

\subsection*{Step 5：人工检查输出}

重点看四件事：

\begin{itemize}
  \item mask 是否只覆盖动态目标；
  \item 细结构是否被漏掉；
  \item 背景是否明显模糊；
  \item 是否有时间闪烁。
\end{itemize}

这个步骤不能省，因为报告里需要 qualitative analysis，不是只贴指标。

\subsection*{Step 6：准备 Part 2}

先运行：

\begin{Verbatim}[fontsize=\small]
python3 scripts/project3.py part2-prepare \
  --config configs/project3.yaml \
  --video data/raw/wild/corridor.mp4 \
  --experiment wild_corridor
\end{Verbatim}

然后再做下面几件事：

\begin{itemize}
  \item clone 或安装 \texttt{SAM 2}；
  \item clone 或安装 \texttt{ProPainter}；
  \item 在 \texttt{configs/project3.yaml} 中写入服务器上的真实路径和 checkpoint；
  \item 根据输出的推荐命令模板真正执行模型。
\end{itemize}

\subsection*{Step 7：计算指标}

一旦有 GT 支持的数据，就运行：

\begin{itemize}
  \item \texttt{eval-mask}；
  \item \texttt{eval-video}。
\end{itemize}

目标是尽快得到：

\begin{itemize}
  \item \texttt{JM / JR}；
  \item \texttt{PSNR / SSIM}。
\end{itemize}

\subsection*{Step 8：最后再启动 Part 3}

不要过早做 diffusion 扩展。

应当满足以下条件后再开始：

\begin{itemize}
  \item Part 1 已稳定；
  \item Part 2 已稳定；
  \item 你们已经能明确指出 ProPainter 的真实 failure case。
\end{itemize}

这样才能保证 Part 3 是一个“有针对性的扩展”，而不是变成无底洞。

\section{两人小组的推荐分工}

\subsection*{同学 A}

\begin{itemize}
  \item 搭服务器环境；
  \item 跑 Part 1 baseline；
  \item 检查 wild / bmx / tennis 的输出效果；
  \item 整理 baseline 的可视化图。
\end{itemize}

\subsection*{同学 B}

\begin{itemize}
  \item 整理数据集目录；
  \item 准备报告模板与 references；
  \item 阅读 SAM 2 / ProPainter 的安装和运行要求；
  \item 准备 progress presentation 内容。
\end{itemize}

等 Part 1 稳定后，两人再一起合流做 Part 2。

\section{最后的实用提醒}

\begin{itemize}
  \item 不要在 Part 1 和 Part 2 还没稳定前就过度投入 Part 3；
  \item 不要把定性图拖到最后再整理；
  \item 不要拖延 references 管理，因为课程明确要求引用给定文献；
  \item 不要默认服务器路径和本地路径一致，迁移后要检查 \texttt{configs/project3.yaml}。
\end{itemize}

\section{结论}

现在这个仓库的定位不是“已经完成所有实验”，而是“已经把项目工程骨架、baseline 逻辑、评测接口和文档模板都搭好”。你们下一阶段的重点不是继续本地开发，而是把这套结构迁移到服务器，先跑通 Part 1，再逐步接入 Part 2 和 Part 3。

\end{document}
